\documentclass[a4paper,14pt]{article}
\usepackage[utf8]{inputenc}
\usepackage[russian]{babel} 
\usepackage{graphicx} 
\usepackage{amsmath}
\usepackage{amssymb}
\usepackage{siunitx}% Для погрешностей
\usepackage{float}
\usepackage[colorlinks=true, linkcolor=blue]{hyperref}
\usepackage{enumitem}
\usepackage{longtable}
%\documentclass[fleqn]{article}

\renewcommand{\floatpagefraction}{0.99}
\renewcommand{\topfraction}{0.99}      


% Подключаем пакет для управления колонтитулами
\usepackage{scrlayer-scrpage}
\pagestyle{scrheadings}

%\clearpagemainheadfoot % Очищаем стандартные колонтитулы
\cfoot{\thepage}
% Настраиваем верхний правый угол (Right Head)
\rofoot{\includegraphics[width=2.5cm]{Mir_insignia.png}}
% Вместо \rohead можно использовать:
% \lehead (левый верхний), \rofoot (правый нижний), \lefoot (левый нижний)



\title{ Фильмы 70/80х}


\author{Батанин Артём Александрович}
\date{\today}


\begin{document}

\maketitle

%\newpage




\section{Героические}
"Есть такая профессия - родину защищать"

Начнем с героических фильмов. В своём периоде я рассмотрел фильм "Офицеры".
Главная идея фильма - показать как через поколения передаётся память о героях прошлого.

--- Герои ---

Главные герои - Алексей Трофимов, его жена Люба, их друг - Иван Варавва.
Сын - Егор Трофимов, Маша Белкина - жена Егора. Внук - Иван Трофимов.


--- Краткий пересказ сюжета ---


Алексей и Иван познакомились в средней азии во время гражданской войны.
Там у Алексея и Любы рождается сын Егор (назвали в честь погибшего офицера)
Варавва на прощание попросил назвать внука Иваном.

Проходит несколько лет, Егор уже вырос и влюбляется в Машу Белкину.
Начинается Великая Отечественная Война - Любу назначают главным санитаром поезда.
На одной из станций к ней подходит Маша, вручает маленького Ваню и заявляет что отправляется на фронт.
На линии фрона поезд подбирает раненых танкистов недалеко от места прорыва Немцев.
Люба подписывала похоронки. На фотографии одного неопознанного солдата она узнала своего сына.

Проходит несколько лет. Алексей командует кантемировской дивизией - ему очень дорога эта дивизия - в ней служил его сын.
Иван учится в суворовском военном училище. Однажды приезжает начальство из Москвы во главе с Иваном Вараввой. его знакомят с Ваней - "у Вараввы перед глазами
проносится весь тот путь, который они прошли от той самой встречи в ближней азии" - он едва сдерживает чувства.

\subsection{Верность долгу и служение Отчизне}
Верность долгу и служение Отчизне демонстрируется в сценах, где Маша Блинова вручает младенца Любе. Она понимала, что её место на фронте.

\subsection{Настоящая дружба}
Дружба Ивана Вараввы и Алексея Трофимова

\subsection{Преемственность поколений}
Было в начале (Основной сюжет фильма)


\subsection{Память о павших}
Егора назвали в честь командира отряда Алексея в средней азии.



\subsection{Стойкость Женщин}
Маша Блинова вручает младенца Любе т.л понимала, что её место на фронте.
Люба находит своего сына во время подписания похоронок.


\section{Лирические}
"Чито-Гврито, Чито-маргалито"

Мимино (Мимино-сокол)

Основная идея фильма -  родина дороже мегаполисов и карьеры.

\subsection{Герои}
Валико - пилот вертолёта в Грузинской деревне
Рубик-джан - водитель Краза
Иван Сергеевич Волохов - Бывший фронтовик. Принял Валико за сына погибшего однополчанина.


\subsection{Сюжет}
Валико работает пилотом самолета в грузинской деревне. Однажды он решает исполнить свою мечту - переквалифицироваться
на дальнемагестрального пилота.
Он пребывает в Москву, где знакомится с Рубик-джаном (водителем красного КРАЗа).
В Москве с ними случаются некоторые приключения, но в итоге Валико действительно становится дальнемагистральным пилотом ИЛ-62.
На одном из рейсов в западный Берлин он зашёл в телефонную будку и попросил позвонить в Телави, но по ошибке его соединяют с
Тель-Авивом. Он поговорил с Мигратом из Грузинской-ССР в Израиль, они спели народную грузинскую песню. Валико понял, что на родине ему будет лучше, после чего вернулся в свою деревню.

\subsection{Ценности}
\begin{itemize}
    \item Любовь к родине и зов предков
    \item Настоящая дружба
    \item Человечность, эмпатия (молодая адвокат борется за Валико в суде)
    \item Внутренняя чистота
\end{itemize}




\section{Комедийные}
Кавказская пленница

\subsection{Герои}

\begin{itemize}
    \item Шурик
    \item Эдик (Будь проклят тот день, когда я сел за баранку этого пылесоса)
    \item Трус, Балбес, Бывалый - Вицин, Никулин, Моргунов
    \item Саахов
\end{itemize}
\subsection{Ценности}

\begin{itemize}
    \item Личная свобода и равноправие женщин
    \item Верность закону и справедливость (Даже старые традиции - не оправдание перед законом)
    \item Настоящая дружба и самоотверженность (Дружба Шурика и Эдика, самоотверженная помощь шурика Нине)
\end{itemize}

\subsection{Цитаты}
«Жить, как говорится, хорошо!» --- «А хорошо жить ещё лучше!»
«Минуточку! Мин-нут-точку! Будьте добры, помедленнее, я з-записую…»
"Птичку жалко"
«Эть, будь проклят тот день, когда я сел за баранку этого пылесоса!»
\end{document}